\chapter{NLP Lanjutan: Di Balik Layar LLM}
\label{chap:nlp_adv}

Di bab awal, kita belajar bahwa LLM adalah "burung beo cerdas". Sekarang, mari kita lihat anatomi otaknya.

\section{Evolusi NLP: RNN vs Transformer}

\subsection{RNN (Recurrent Neural Networks)}
Model lama membaca kata satu per satu secara berurutan.
\textit{"Saya"} $\rightarrow$ \textit{"suka"} $\rightarrow$ \textit{"makan"}.
Kelemahan: Lupa konteks awal jika kalimat terlalu panjang.

\subsection{Transformer (Attention Mechanism)}
Terobosan Google (2017). Transformer membaca seluruh kalimat sekaligus dan memberikan "perhatian" (attention) pada hubungan antar kata, tidak peduli seberapa jauh jaraknya.

\begin{quote}
"Attention is All You Need."
\end{quote}

\section{BERT vs GPT}

\begin{itemize}
    \item **BERT (Bidirectional Encoder):** Membaca dari kiri-ke-kanan DAN kanan-ke-kiri. Bagus untuk \textit{memahami} teks (klasifikasi, analisis sentimen).
    \item **GPT (Generative Pre-trained Transformer):** Hanya membaca kiri-ke-kanan (Decoder). Bagus untuk \textit{membuat} teks (generasi).
\end{itemize}

\section{Fine-Tuning vs RAG}

Kapan harus melatih ulang model?

\begin{center}
\begin{tabular}{|l|l|l|}
\hline
\textbf{Fitur} & \textbf{Fine-Tuning} & \textbf{RAG} \\
\hline
\textbf{Tujuan} & Mengubah gaya/perilaku & Menambah pengetahuan \\
\hline
\textbf{Biaya} & Mahal & Murah \\
\hline
\textbf{Update Data} & Butuh training ulang & Instan (update database) \\
\hline
\textbf{Contoh} & Chatbot medis & Chatbot Customer Service \\
\hline
\end{tabular}
\end{center}

\begin{tipbox}
Gunakan RAG untuk fakta. Gunakan Fine-Tuning untuk gaya bahasa.
\end{tipbox}
